\section{RNN 架构扩展}

\subsection{梯度问题不仅限于 RNN}

\begin{figure}[H]
\centering
\includegraphics[width=0.95\textwidth]{slides-images/slide-025.jpg}
\caption{梯度消失/爆炸不仅是 RNN 的问题：所有深层网络都面临类似挑战。残差连接（ResNet）是深层前馈网络的解决方案\protect\footnotemark}
\end{figure}
\footnotetext{Slides 第26页。}

梯度消失和爆炸并非 RNN 独有的问题。在任何深层网络中，梯度都需要经过多层传递，同样可能出现指数级的衰减或增长。这也是为什么早期深层神经网络难以训练的原因。

\begin{knowledgebox}{残差连接：深层网络的类似解决方案}
\textbf{残差网络}（ResNet, He et al., 2015）通过在层间添加\textbf{恒等映射}（skip connection）来解决深层前馈网络的梯度问题：$\mathcal{F}(x) + x$。这与 LSTM 的思路异曲同工——都是通过\textbf{加法}提供信息直通的路径。相关的变体包括：
\begin{itemize}
    \item \textbf{DenseNet}：每一层与所有后续层都直接连接
    \item \textbf{Highway Network}（Srivastava et al.，Schmidhuber 团队）：在残差连接上增加门控，类似 LSTM 的门控思路
\end{itemize}
\end{knowledgebox}

\subsection{RNN 的多种应用}

一旦有了 RNN（尤其是 LSTM），它就可以用于各种序列相关的 NLP 任务：

\begin{figure}[H]
\centering
\includegraphics[width=0.95\textwidth]{slides-images/slide-029.jpg}
\caption{RNN 的多种应用场景：标注任务、句子编码、生成任务等\protect\footnotemark}
\end{figure}
\footnotetext{Slides 第30页。}

\begin{itemize}
    \item \textbf{序列标注}：词性标注（POS Tagging）、命名实体识别（NER）——在每个位置输出一个标签
    \item \textbf{句子编码}：情感分类等——将整个句子编码为一个固定长度的向量，用于分类
    \item \textbf{条件生成}：机器翻译、语音识别、文本摘要——在给定输入条件下生成输出序列
\end{itemize}

\subsection{句子编码：RNN 作为句子表示}

\begin{figure}[H]
\centering
\includegraphics[width=0.95\textwidth]{slides-images/slide-030.jpg}
\caption{RNN 用于情感分类：将整个句子编码为一个向量，用于判断正面/负面情感。实践中取所有隐藏状态的均值或逐元素最大值通常优于仅用最后一步\protect\footnotemark}
\end{figure}
\footnotetext{Slides 第31页。}

对于句子级别的分类任务（如情感分析），可以用 RNN 处理整个句子，然后使用得到的表示进行分类。最简单的方法是取最后一步的隐藏状态 $h^{(T)}$，但实践中更好的做法是对所有时间步的隐藏状态取\textbf{均值}或\textbf{逐元素最大值}。

\subsection{双向 RNN}

单向 RNN 在每个位置只能看到左侧的上下文。但对于很多任务（如序列标注），我们希望每个位置的表示能综合\textbf{左右两侧}的信息。

\begin{figure}[H]
\centering
\includegraphics[width=0.95\textwidth]{slides-images/slide-034.jpg}
\caption{双向 RNN：分别运行前向和后向 RNN，在每个位置拼接两个方向的隐藏状态\protect\footnotemark}
\end{figure}
\footnotetext{Slides 第35页。}

\begin{figure}[H]
\centering
\includegraphics[width=0.95\textwidth]{slides-images/slide-035.jpg}
\caption{双向 RNN 的简化图示：双向箭头表示每个位置的隐藏状态是前向和后向状态的拼接\protect\footnotemark}
\end{figure}
\footnotetext{Slides 第36页。}

\begin{importantbox}{双向 RNN（Bidirectional RNN）}
\begin{itemize}
    \item 分别运行一个\textbf{前向} LSTM（从左到右）和一个\textbf{后向} LSTM（从右到左）
    \item 在每个位置，将两个方向的隐藏状态\textbf{拼接}：$h^{(t)} = [\overrightarrow{h}^{(t)}; \overleftarrow{h}^{(t)}]$
    \item 适用于\textbf{编码}任务（分类、标注等），但\textbf{不能}用于语言生成（因为后向 RNN 需要看到未来的词）
\end{itemize}
\end{importantbox}

\subsection{多层（堆叠）RNN}

\begin{figure}[H]
\centering
\includegraphics[width=0.95\textwidth]{slides-images/slide-037.jpg}
\caption{多层 RNN（Stacked RNN）：多层 LSTM 叠加，下层的隐藏状态作为上层的输入，实现更强的特征提取\protect\footnotemark}
\end{figure}
\footnotetext{Slides 第38页。}

将多层 RNN 堆叠可以提升模型的表示能力，类似于深层前馈网络通过多层实现更高级的特征抽象。

\begin{warningbox}{RNN 时代的``深度''有限}
与现代 Transformer 动辄几十甚至上百层不同，RNN 时代人们通常只使用 2--3 层。从 1 层到 2 层的提升很明显，但继续增加层数的收益递减。Manning 指出，这一状况在 Transformer 时代已完全改变——现代语言模型使用非常深的 Transformer 网络。
\end{warningbox}

\subsection{本章小结}

\begin{itemize}
    \item 梯度问题不仅限于 RNN，深层前馈网络的解决方案（残差连接）与 LSTM 的加法更新思路一脉相承
    \item \textbf{双向 RNN} 提供了每个位置的完整上下文表示，适合编码任务但不能用于生成
    \item \textbf{多层 RNN} 通过堆叠增强表示能力，RNN 时代通常用 2--3 层
    \item LSTM 在 2013--2015 年间成为几乎所有 NLP 任务的主导方法
\end{itemize}

%% ========== 机器翻译 ==========

